% =====================================================================
% Слайд 6 — Baseline 1: OLTW (наработка лаборатории)
% =====================================================================
\begin{frame}{Baseline 1: OLTW \small(Dixon 2005)}
  \small\setlength{\abovedisplayskip}{2pt}
  \setlength{\belowdisplayskip}{2pt}

  \textbf{On-Line Time Warping} --- инкрементальный вариант DTW.

  \begin{columns}[T,onlytextwidth]
    \column{0.52\textwidth}
    \begin{itemize}\setlength{\itemsep}{0.25em}
      \item Локальная стоимость:
            $\displaystyle c(i,t) = |\,p_i - o_t|$.
      \item Рекурсия:\\
            $\displaystyle D[i,t] = c(i,t) +
             \min\bigl(D[i{-}1,t{-}1],\, D[i{-}1,t],\, D[i,t{-}1]\bigr)$.
      \item Предсказание: $\hat\imath_t = \argmin_i D[i,t]$.
      \item Память $O(N)$ --- хранятся только две колонки.
      \item Монотонизация $+$ \emph{RunCount} fail-safe: если позиция
            $k$ событий не двигается, форсируем шаг вперёд.
    \end{itemize}

    \column{0.48\textwidth}
    \centering
    % СХЕМА 3 — OLTW dynamic programming
\begin{tikzpicture}[
  every node/.style={font=\footnotesize},
  cell/.style={draw, thick, minimum width=1.1cm, minimum height=0.7cm,
               inner sep=1pt, fill=white},
  arr/.style={-{Latex[length=2mm]}},
]
  % Header
  \node[font=\scriptsize] at (0, 3.0) {$t{-}1$};
  \node[font=\scriptsize] at (1.6, 3.0) {$t$};

  % Grid: rows i = 0..4
  \foreach \r/\i in {0/4, 1/3, 2/2, 3/1, 4/0} {
    \node[cell] (p\i) at (0, 2.5-\r*0.75) {$D[\i,t{-}1]$};
    \node[cell] (c\i) at (1.6, 2.5-\r*0.75) {$D[\i,t]$};
    \node[font=\scriptsize] at (-0.95, 2.5-\r*0.75) {$i{=}\i$};
  }

  % Arrows into c2 (D[2,t])
  \draw[arr] (p1) -- (c2);     % diagonal D[i-1,t-1]
  \draw[arr] (p2) -- (c2);     % horizontal D[i,t-1]
  \draw[arr] (c1) -- (c2);     % vertical D[i-1,t]
  \node[font=\scriptsize, right=4pt of p1, anchor=west] {\itshape diag};
  \node[font=\scriptsize, right=4pt of p2, anchor=west] {\itshape horiz};
  \node[font=\scriptsize, right=4pt of c1, anchor=west, xshift=8pt] {\itshape vert};

  % Highlight target cell
  \draw[very thick] (c2.south west) rectangle (c2.north east);

  % Equation below
  \node[align=left, font=\scriptsize] at (1.0, -1.0)
       {$D[i,t] \;=\; c(i,t) \;+\; \min\bigl(D[i{-}1,t{-}1],\; D[i{-}1,t],\; D[i,t{-}1]\bigr)$};
  \node[align=center, font=\scriptsize\itshape] at (1.0, -1.55)
       {Predicted index $= \arg\min_i D[i,t]$};
\end{tikzpicture}

    \figcaption{Predicted index $=$ argmin по текущему столбцу.}
  \end{columns}

  \vspace{0.3em}
  {\footnotesize \textbf{Плюсы:} простота, устойчивость к rubato,
   всегда выдаёт ответ.\quad
   \textbf{Минусы:} нет confidence, плохо ловит структурные ошибки.}
\end{frame}

